\chapter{Ketogenic Diet}

\section*{Definition and Overview}
The ketogenic diet is a very low-carbohydrate, high-fat diet that forces the body to burn fat for fuel instead of glucose, producing ketones. It has a well-established medical role in the treatment of drug-resistant epilepsy in children, where it can dramatically reduce seizures, and it is sometimes used in adults with epilepsy. In recent years it has become popular for weight loss, and it produces rapid initial weight loss, although evidence on long-term benefits is mixed and the diet is difficult to sustain. Because it is restrictive and can cause nutritional deficiencies and other effects, it should be followed under medical and dietetic supervision, especially in children.

\section*{Epidemiology}
The ketogenic diet was developed in the 1920s for epilepsy and remains an important treatment for the 20 to 30 per cent of children with epilepsy whose seizures are not controlled by medicines. Its use declined with the arrival of anticonvulsant drugs but has increased again, with modified versions such as the modified Atkins diet being easier to follow. As a weight-loss approach, low-carbohydrate and ketogenic diets are among the most popular in Australia and worldwide, and many people use them to lose weight, with reported benefits in blood sugar and triglycerides.

\section*{Aetiology and Pathophysiology}
Normally the body uses glucose from carbohydrates as its main fuel, with glucose stored as glycogen for short-term needs. When carbohydrate intake is very low, glycogen runs out within a day or two, and the body shifts to burning fat, producing ketones in the liver. In epilepsy, ketones and the metabolic changes of the diet are thought to reduce nerve excitability and seizure frequency, though the exact mechanisms are not fully understood. In weight loss, the diet suppresses appetite and causes rapid initial loss of water and glycogen, followed by fat loss, and the overall energy deficit drives the weight change. Ketosis is a normal metabolic state, but the diet is restrictive and can be hard to sustain, and it carries risks if not followed properly.

\section*{Clinical Features}
\begin{itemize}
    \item Very low carbohydrate intake, typically under 50 grams daily, with high fat and moderate protein
    \item Weight loss, often rapid in the first weeks
    \item Reduced appetite and a feeling of fullness
    \item Keto flu in the first days: headache, fatigue, nausea, and brain fog
    \item Bad breath and changes in body odour from ketones
    \item Reduced seizures in many children with drug-resistant epilepsy
    \item Constipation, and in some people, kidney stones
    \item Improved blood sugar and triglycerides in some people
    \item Difficulty sustaining the diet long term
\end{itemize}

\section*{Differential Diagnosis}
The ketogenic diet must be distinguished from other eating approaches and assessed for suitability.
\begin{itemize}
    \item \textbf{Low-carbohydrate diets:} moderate restriction, easier to sustain
    \item \textbf{Mediterranean and balanced diets:} recommended for general health
    \item \textbf{Weight-loss programs:} structured calorie control
    \item \textbf{Medical treatment of epilepsy:} medicines and other options
    \item \textbf{Eating disorders:} restrictive diets can trigger or worsen them
    \item \textbf{Medical conditions:} diabetes, kidney disease, and pregnancy require caution
\end{itemize}

\section*{When to Seek Medical Care}
The ketogenic diet for epilepsy is always started under specialist medical and dietetic supervision, and families must not begin it without this. People considering it for weight loss should discuss it with a doctor, especially those with diabetes, kidney disease, heart disease, or a history of eating disorders, and it is not suitable in pregnancy or breastfeeding. Anyone on the diet who develops concerning symptoms should seek review.

\textbf{Call 000 (or local emergency services) for:}
\begin{itemize}
    \item Confusion, collapse, or unusual drowsiness
    \item Severe abdominal pain or vomiting
    \item Difficulty breathing or chest pain
    \item Signs of very low blood sugar in someone with diabetes
    \item Seizures or convulsions in a child on the diet
\end{itemize}

\section*{Diagnosis and Investigations}
\begin{itemize}
    \item \textbf{Medical and dietetic assessment:} suitability, goals, and monitoring plan
    \item \textbf{Blood tests:} glucose, electrolytes, kidney and liver function, and lipids
    \item \textbf{Ketone monitoring:} blood or urine ketones to confirm ketosis
    \item \textbf{Growth monitoring in children:} on the epilepsy diet
    \item \textbf{Review of medicines:} interactions and adjustment, especially for diabetes
    \item \textbf{Screening for eating disorders:} before restrictive diets
    \item \textbf{Ongoing review:} of weight, nutrition, and side effects
\end{itemize}

\section*{Management}
\begin{itemize}
    \item \textbf{Medical epilepsy treatment:} the diet is started in hospital with dietetic support and close monitoring
    \item \textbf{Dietitian guidance:} structuring meals, supplements, and micronutrient monitoring
    \item \textbf{Modified versions:} the modified Atkins diet and MCT diet, which are less restrictive
    \item \textbf{Supplements:} vitamins, minerals, and electrolytes as needed
    \item \textbf{Managing side effects:} fluids, fibre, and treatment of constipation and stones
    \item \textbf{Weight-loss use:} short-term with dietetic review and a plan for maintenance
    \item \textbf{Stopping the diet:} gradual reintroduction of carbohydrates under guidance
    \item \textbf{Monitoring:} regular review of health and, in epilepsy, seizure control
\end{itemize}

\section*{Complications}
\begin{itemize}
    \item Nutritional deficiencies if the diet is not planned properly
    \item Constipation, kidney stones, and dehydration
    \item Raised cholesterol in some people
    \item Growth problems in children if not monitored
    \item Bone density changes over time
    \item Difficulties sustaining the diet and regaining weight
    \item Risks in people with diabetes, kidney disease, or eating disorders
    \item Ketoacidosis, which is different from nutritional ketosis and dangerous
\end{itemize}

\section*{Prognosis and Outcomes}
For children with drug-resistant epilepsy, the ketogenic diet is highly effective in a substantial proportion, with around half experiencing at least a 50 per cent reduction in seizures, and it can be life-changing. For weight loss, it produces rapid initial results, but long-term weight loss is similar to other diets when adherence is considered, and most people find it difficult to sustain. With proper supervision, the diet is safe for its medical uses, and it can be followed for weight loss when individualised and monitored.

\section*{Prevention and Self-Care}
\begin{itemize}
    \item Use the ketogenic diet for epilepsy only under specialist supervision
    \item Seek dietetic advice before starting for weight loss
    \item Discuss it with your doctor if you have any medical condition
    \item Do not use the diet in pregnancy or breastfeeding
    \item Stay well hydrated and manage electrolytes
    \item Include non-starchy vegetables and fibre where the diet allows
    \item Take supplements as advised
    \item Monitor your health with regular blood tests
    \item Plan how you will maintain weight after stopping
    \item Seek help early for side effects or difficulties
\end{itemize}

\section*{Sources and Further Reading}
The following reputable sources provide further information for healthcare students and patients seeking reliable, current advice about the ketogenic diet.
\begin{itemize}
    \item Epilepsy Action Australia. \textit{Ketogenic diet for epilepsy.} https://www.epilepsy.org.au/
    \item Ketogenic Diet Foundation (or Kids with Epilepsy). \textit{Ketogenic diet information.} https://www.epilepsysmart.com.au/
    \item Dietitians Australia. \textit{Ketogenic diet advice.} https://dietitiansaustralia.org.au/
    \item Healthdirect Australia. \textit{Ketogenic diet.} https://www.healthdirect.gov.au/
    \item NHS. \textit{Ketogenic diet.} https://www.nhs.uk/
    \item Charlie Foundation. \textit{Ketogenic diet for epilepsy.} https://charliefoundation.org/
\end{itemize}
